\newcommand{\LavaTitle}{Class-Asymmetric Robustness and Transformation-Order Effects in Lightweight Audio Deepfake Detection}
\newcommand{\LavaAbstract}{Speech deepfake detectors encounter noise and lossy coding in succession. We examine how channel augmentation changes ranking and classwise decisions under ordered processing. Paired MobileNetV3Small--LSTM detectors trained on 640 ASVspoof 2019 LA recordings shared initialization, source order, optimization budget, and the clean-development selection rule within each of two seed blocks. A locked within-release evaluation used 1,088 recordings from speakers absent from the training and development cohorts, four conditions, and four development-selected models. Mean stressed ROC-AUC increased from 0.8308 to 0.8875 and from 0.7498 to 0.8948. Stressed FAKE misses fell from 1,370 to 587 and from 815 to 453 out of 2,496 transformed views per block. Clean FAKE misses increased and clean ROC-AUC declined. Descriptive subgroup analysis found increased stressed misses on A17--A19 in both blocks. In the 9207 noise-before-Opus condition, REAL false alarms fell from 226 to 89 of 256 while FAKE misses rose from 38 to 93 of 832. These results support evaluating ranking and classwise errors together at the intended operating point and testing processing order on matched recordings.}
\newcommand{\LavaKeywords}{Audio forensics, channel augmentation, codec robustness, spoofing countermeasures}
